\section{Beurteilung des Kursverlaufs}
\label{sec:kurs}

\margentabelle

Der Kurs folgt bei Neste nicht dem Umsatz und nicht dem Ergebnis des Jahres, sondern zwei
Gr\"o{\ss}en zugleich: der Verkaufsmarge von Renewable Products und dem Multiple, das der
Markt auf das erwartete Wachstum zahlt. Tabelle~\ref{tab:marge} und Tabelle~\ref{tab:ev}
trennen beide.

\textbf{2020--2023: Das Multiple fiel, die Marge stieg.} Im Jahr 2020 stieg der Kurs um
\KursDeltaXX{} bei einer um \MargeDeltaXX{} ver\"anderten Marge; das Verh\"altnis von
Unternehmenswert zu EBITDA erreichte \EvMax{}$\times$ (Tabelle~\ref{tab:ev}). Bis 2023 stieg
die Marge in jedem Jahr, das EBITDA 2022 und 2023 auf den H\"ochststand der Reihe, und der Kurs
fiel dennoch um \KursDeltaXXI{}, \KursDeltaXXII{} und \KursDeltaXXIII{}. Was fiel, war das
Multiple, bis auf \EvMin{}$\times$ im Jahr 2023. \bk{Der Markt hatte 2020 ein Wachstum
bezahlt, das sich in den Ergebnissen dann zeigte, aber nicht in dem Umfang, den der Preis
vorwegnahm.}

\textbf{2024: Der Indikator fiel.} Die Marge brach um \MargeDeltaXXIV{} ein, der Kurs um
\KursDeltaXXIV. Das ist das einzige Jahr der Reihe, in dem der F\"uhrungsindikator selbst den
Kurs erkl\"art. Neste hatte die Margenspanne im Jahresverlauf zweimal gesenkt\Q{GBXXIV}{98}.

\textbf{2025 bis heute: Erst das Multiple, dann die Marge.} 2025 stieg die Marge um nur
\MargeDeltaXXV, der Kurs um \KursDeltaXXV; das Multiple stieg von \EvXXIII{}$\times$ (2023) auf \EvXXV{}$\times$ zum
Jahresende 2025. Im ersten Halbjahr 2026 folgte die
Marge mit \je{\HjRpMarge}{USD/t}, und der Kurs stieg seit Jahresende um \KursSeitJahresende.
Auf das EBITDA der letzten zw\"olf Monate gerechnet liegt das Multiple heute bei
\EvHeuteLtm{}$\times$, auf Rang \Pz{\EvRangLtm} der eigenen Reihe.

\textbf{Die Reaktion auf den Rekord.} Am Tag des Halbjahresberichts (24.07.2026), der das
h\"ochste Quartalsergebnis der Unternehmensgeschichte meldete\Q{HJ}{1}, fiel die Aktie von
\je{\KursReaktionVor}{EUR} auf \je{\KursReaktionNach}{EUR}, um \ReaktionHj{}
(Yahoo Finance). \bk{Der Markt hat den Rekord nicht fortgeschrieben, sondern als
H\"ohepunkt gelesen. Dieselbe Lesart zeigt der Konsens mit einem f\"ur 2027 um
\EpsRueckgangXXVII{} niedrigeren Ergebnis je Aktie. Der Kurs ist seither dennoch weiter
gestiegen, begleitet von angehobenen Kurszielen mehrerer H\"auser (Abschnitt~1.4).}
